% Generated draft: requires mathematical and source review.
\documentclass[11pt]{article}
\usepackage[margin=25mm]{geometry}
\usepackage{fontspec}
\setmainfont{Latin Modern Roman}
\usepackage{amsmath,amssymb,mathtools,unicode-math}
\setmathfont{Latin Modern Math}
\usepackage{enumitem,needspace,xurl,hyperref}
\setlength{\parindent}{0pt}
\setlength{\parskip}{5pt}
\setlength{\emergencystretch}{3em}
\newcommand{\N}{\mathbb{N}}
\newcommand{\Z}{\mathbb{Z}}
\newcommand{\Q}{\mathbb{Q}}
\newcommand{\R}{\mathbb{R}}
\newcommand{\C}{\mathbb{C}}
\newcommand{\F}{\mathbb{F}}
\newcommand{\A}{\mathbb{A}}
\newcommand{\PP}{\mathbb{P}}
\newcommand{\E}{\mathbb{E}}
\newcommand{\eps}{\varepsilon}
\newcommand{\dd}{\,\mathrm{d}}
\begin{document}
\section*{UnsolvedMath ID: 30003831}
\textbf{Source problem code:} catalogue entry
\subsection{Definitions and mathematical statement}
Let \(k\) be a field, let \(a_1,\ldots,a_n\) be a finite list of elements algebraically independent over \(k\), and put \(K=k(a_1,\ldots,a_n)\). Algebraic independence means that no nonzero polynomial in \(k[X_1,\ldots,X_n]\) vanishes at \((a_1,\ldots,a_n)\). Let \(L/K\) be an explicitly given finite separable extension. Fix a separable closure \(K^{\mathrm{sep}}\) of \(K\) and a \(K\)-embedding of \(L\) into it. The Galois closure \(\widetilde L\) is the subfield of \(K^{\mathrm{sep}}\) generated over \(K\) by the images of all \(K\)-embeddings \(L\hookrightarrow K^{\mathrm{sep}}\). It is a finite Galois extension of \(K\). Its Galois group is
\[
G=\operatorname{Gal}(\widetilde L/K)
 =\{\sigma\in\operatorname{Aut}(\widetilde L):\sigma|_K=\operatorname{id}_K\}.
\]
The isomorphism class of \(G\) does not depend on the choices above.

\textbf{Problem as supplied.} Construct an algorithm which, for every field \(k\), every such finite list \(a_1,\ldots,a_n\), and every explicitly given finite separable extension \(L/K\), computes \(\operatorname{Gal}(\widetilde L/K)\).

\textbf{Unresolved algorithmic qualifications.} The source does not specify an effective presentation of the arbitrary field \(k\), available field operations or decision procedures, the encoding of \(L/K\), or the required representation of the output group. Consequently, the displayed request identifies the mathematical target but does not yet specify a precise computational problem. In particular, no computability assumption on \(k\), oracle access, restriction on its characteristic, or meaning of ``explicitly given'' is added here. These choices require clarification before a uniform algorithmic assertion can be formulated.
\subsection{Short English statement}
Given a finite separable extension of a rational function field over an arbitrary field, compute the automorphism group of its Galois closure over that rational function field. The supplied formulation leaves the computational access to the base field, input encoding, and output format unspecified. Its open status is supplied dated curation, not an independently verified conclusion.
\subsection{Sources}
\begin{itemize}
\item Local UnsolvedMath registry, supplied complete source record, permanent ID \(30003831\), exact source problem code \texttt{catalogue entry}; title: \emph{Computing Galois Closures over Rational Function Fields}. Collection: \emph{Oberwolfach Reports Open Problems}. No individual proposer is supplied; the recorded proposed year is 2018.
\item Supplied report citation: \emph{Field Arithmetic (2019)}, Oberwolfach Reports, DOI 10.4171/owr/2018/25, pp. 1517--1577. \url{https://doi.org/10.4171/owr/2018/25}. The supplied background describes these reports as documenting mathematical workshops and their open-problem discussions.
\item Supplied literature source, labelled ``Oberwolfach Report 2019'': \url{https://ems.press/content/serial-article-files/46747}.
\item Supplied literature source, labelled ``Entin'': \url{https://arxiv.org/abs/1805.05454}. No title or specific supporting result is supplied.
\item Literature assessment dated exactly \texttt{2026-08-21}: supplied dated curation, not independently verified here. It labels the problem open, describes the excerpt as a general algorithmic question truncated before a claim or resolution, and states that existing work treats special polynomial and geometric families rather than the full problem. The background explicitly requests independent verification of time-sensitive status and cited claims.
\item The registry labels the statement \texttt{corrected\_verified} and records that the official report was checked directly, with definitions restored and source-dependent wording removed. This is a supplied source-verification assertion; no independent source visit or verification is claimed here.
\end{itemize}
\end{document}
